% Image credits for the photographs and AI illustrations of Book 5
% (University Physics, Year 3). Machine-readable license records live in
% images/book5/CREDITS.md; keep the two files in sync.
\chapter*{图片来源}

本书中的插图均为原创。照片依据各自的自由许可使用，谨向其作者致谢：

\begin{itemize}
  \item \emph{先贤祠中的傅科摆}（第 1 章）：
        Olga Khomitsevich，Wikimedia Commons，CC~BY~2.0。
  \item \emph{阿尔伯特·爱因斯坦}（第 4 章）：Orren Jack Turner，
        1947，公有领域。
  \item \emph{正电子的发现}（第 5 章）：Carl~D.
        Anderson，1932，公有领域。
  \item \emph{1927 年索尔维会议}（第 8 章）：Benjamin
        Couprie，公有领域。
  \item \emph{第一个玻色–爱因斯坦凝聚体}（第 19 章）：
        NIST/JILA/CU-Boulder，公有领域。
  \item \emph{宇宙微波背景}（第 20 章）：
        NASA / WMAP Science Team，公有领域。
  \item \emph{第一个晶体管的复制品}（第 24 章）：
        贝尔实验室展品，公有领域。
  \item \emph{玛丽·居里}（第 25 章）：Henri Manuel，
        c.~1920，公有领域。
  \item \emph{建造中的 CMS 探测器}（第 26 章）：
        Julian Williams，Wikimedia Commons，版权所有者允许自由使用。
  \item \emph{蟹状星云}（第 27 章）：NASA，ESA，
        J.~Hester 和 A.~Loll（Arizona State University），公有
        领域。
\end{itemize}

每张照片的完整来源链接与许可证 URL 均列于本书代码库中的
\texttt{images/book5/CREDITS.md}。

\bigskip
除照片外，第 1--8 章和第 11--27 章中共有二十六幅图为 AI 生成的插图，
由 OpenAI 图像生成功能根据本书编辑撰写的提示词制作而成，并在收录前
核查了物理上的准确性。每幅生成图像的完整提示词均记录于代码库中的
\texttt{images/book5/ai/PROMPTS.md}。
\clearpage
